%
\subsection*{Estructura lineal y condiciones}
\label{CH04-B01}
En un intervalo donde el coeficiente principal no se anula y los
coeficientes normalizados son continuos, una EDO lineal de orden $n$
tiene solución general $y=y_h+y_p$. Las soluciones homogéneas forman un
espacio de dimensión $n$. Para segundo orden,
$W(y_1,y_2)=y_1y_2'-y_1'y_2$; si ambas resuelven la homogénea, un
wronskiano no nulo en un punto del intervalo prueba independencia.
Las $n$ condiciones iniciales en un mismo punto fijan una única solución.

Las condiciones en extremos distintos forman un problema de contorno y
pueden dar ninguna, una o varias soluciones. Al imponerlas se obtiene un
sistema algebraico cuya compatibilidad debe comprobarse. Un cero en una
fórmula auxiliar no prueba que la ecuación sea singular: revísense sus
coeficientes antes de concluir. Bastan derivación y sistemas lineales
elementales para estas comprobaciones.

%
\subsection*{Raíces características y bases reales}
\label{CH04-B02}
Para coeficientes constantes, se sustituye $y=e^{rt}$ y se resuelve el
polinomio característico. Una raíz real $r$ de multiplicidad $m$ aporta
$e^{rt},te^{rt},\ldots,t^{m-1}e^{rt}$. Un par
$\alpha\pm i\beta$, $\beta>0$, de multiplicidad $m$ aporta
\[
 t^j e^{\alpha t}\cos(\beta t),\quad
 t^j e^{\alpha t}\sin(\beta t),\qquad j=0,\ldots,m-1.
\]
Estas expresiones se obtienen de
$e^{i\theta}=\cos\theta+i\sin\theta$ y permiten trabajar con constantes
reales. Derívese la combinación antes de imponer los datos y compruébese
el PVI completo. No se divida por un parámetro sin tratar sus posibles
ceros; una raíz repetida exige todas las potencias indicadas.
Consulta: solución de \exref{M2-E0015}.

%
\subsection*{Reducción de orden y dominio}
\label{CH04-B03}
Si una ecuación no contiene $y$, poner $v=y'$ reduce el orden en una
unidad; al recuperar $y$ aparece una constante de integración adicional.
Si se conoce una solución $y_1$ de
$y''+p y'+q y=0$, el cambio $y=y_1v$, en un intervalo donde $y_1\ne0$,
da
\[
 y_1v''+(2y_1'+p y_1)v'=0.
\]
La sustitución $w=v'$ deja una ecuación de primer orden. La regla del
producto explica qué términos se cancelan; no basta memorizar el cambio.

La reducción es local: que una expresión intermedia se divida por
$y_1$ no decide si la solución final puede atravesar un cero de $y_1$.
Una posible prolongación se comprueba en la ecuación original, con la
regularidad y los datos correspondientes. Consulta: solución de
\exref{M2-E0224}.

%
\subsection*{Cómo construir una particular}
\label{CH04-B04}
Con coeficientes constantes, un forzamiento
$e^{\alpha t}[P(t)\cos\beta t+Q(t)\sin\beta t]$ admite un ensayo del
mismo tipo, con ambos polinomios de grado suficiente. Multiplíquese por
$t^m$ si $\alpha+i\beta$ es raíz característica de multiplicidad $m$
($\beta=0$ incluye exponenciales y polinomios). Un anulador puede ampliar
la ecuación; las constantes de la particular se fijan volviendo a la
ecuación original.

Para $y''+p y'+q y=r$, con una base homogénea $y_1,y_2$ y $W\ne0$,
una particular es
\[
 y_p=-y_1(t)\int_{t_*}^{t}\frac{y_2(s)r(s)}{W(s)}\,ds
      +y_2(t)\int_{t_*}^{t}\frac{y_1(s)r(s)}{W(s)}\,ds.
\]
Normalícese primero el coeficiente de $y''$. Las cuadraturas definidas son
válidas aunque no sean elementales. El fallo de un ensayo en un parámetro
no prueba inexistencia. Consultas: soluciones de \exref{M2-E0016} y
\exref{M2-E0303}.

%
\subsection*{Recuperar estados, datos y escalas}
\label{CH04-B05}
Al eliminar una variable de un balance acoplado, consérvese la relación
que permite reconstruirla. El valor inicial de la derivada de la variable
retenida se obtiene del balance original, no se elige libremente. Tras
resolver, reconstruya el estado eliminado con las mismas constantes y
compruebe ambas ecuaciones. Las condiciones de los extremos se traducen
del mismo modo y se revisa si son independientes.

En familias con un parámetro, distínganse un tiempo fijo, un tiempo que
varía con ese parámetro y una garantía para todos los tiempos.
Una cota uniforme $|y_\varepsilon-y|\le B(\varepsilon)$ exige controlar
todo el intervalo considerado. Aproximar el estado no garantiza aproximar
su derivada; ésta requiere su propia comprobación. Son útiles los balances
y las desigualdades elementales. Consultas: soluciones de
\exref{M2-E0313} y \exref{M2-E0334}.
